% --- chapters/gauge_sector.tex --------------------------------------------
\section{Gauge Sector from the \texorpdfstring{$E_8$}{E8} Embedding}
\label{sec:gauge-from-e8}

We realize the SM gauge algebra as
\[
SU(3)_C \times SU(2)_L \times U(1)_Y \;\subset\; SU(5)\subset SO(10)\subset E_8
\]
within the same \emph{24‑cell (base)} $\times$ \emph{fiber Cartan} slice used elsewhere.
Base coordinates carry dynamics; fiber coordinates carry charges via inner
products with Cartan generators.

\paragraph{Hypercharge.}
Inside $SU(5)$ choose
\[
Y \;=\; \mathrm{diag}\!\left(-\tfrac13,-\tfrac13,-\tfrac13,+\tfrac12,+\tfrac12\right)
\quad\text{on the } \mathbf{5} ,
\]
with GUT normalization
\[
g_1=\sqrt{\tfrac{5}{3}}\,g_Y, \qquad
\mathrm{tr}(T^aT^b)=\tfrac12\delta^{ab},\; \mathrm{tr}(Y^2)=\tfrac12\!\cdot\!\tfrac{5}{3}.
\]
Equivalently, along the Pati–Salam chain
\[
Y \;=\; T^3_{R} + \tfrac12(B\!-\!L).
\]

\paragraph{Constructive extraction (fiber‑linear).}
Fix an $E_8$ simple‑root basis and an orthogonal embedding matrix
$P:\mathbb R^8\!\to\!\mathbb R^8$ for
$E_8\to SO(10)\times SU(4)\to SU(5)\times U(1)_\chi$.
For any weight $w\in\Lambda_{E_8}$,
\[
w_{SU(5)} = P\,w, \qquad
Y(w) = \big\langle P^\top Y_{\rm Cartan},\, w\big\rangle_{E_8}.
\]
$SU(3)_C$ and $SU(2)_L$ Dynkin labels are read from the projected components.
In our geometry this is “drop base, keep fiber, apply $P$”.

\paragraph{One family (worked).}
From $SO(10)$, the chirality is a $\mathbf{16}\to \mathbf{10}\oplus\overline{\mathbf{5}}\oplus \mathbf{1}$.
Decomposed to SM:
\[
\mathbf{10}\to (3,2)_{+1/6}\oplus(\overline3,1)_{-2/3}\oplus(1,1)_{+1},\quad
\overline{\mathbf{5}}\to (\overline3,1)_{+1/3}\oplus(1,2)_{-1/2},\quad
\mathbf{1}\to (1,1)_0.
\]
These are exactly
\[
Q:(3,2)_{+1/6},\quad
u^c:(\overline3,1)_{-2/3},\quad
e^c:(1,1)_{+1},\quad
d^c:(\overline3,1)_{+1/3},\quad
L:(1,2)_{-1/2},\quad
(\nu^c):(1,1)_0 .
\]

\paragraph{Anomaly checks (per generation).}
\begin{align*}
&SU(3)^2U(1)_Y: && 3\!\times\!2\cdot\tfrac16 \;+\; 3\cdot(-\tfrac23) \;+\; 3\cdot(+\tfrac13) \;=\;0,\\
&SU(2)^2U(1)_Y: && 3\cdot\tfrac16 \;+\; 1\cdot(-\tfrac12) \;=\;0,\\
&U(1)^3_Y: && \sum_{\rm fam} Y^3 = 0,\qquad
\mathrm{Grav}^2\!-\!U(1)_Y:\ \sum Y = 0,\\
&\text{Witten global }SU(2): && \#\text{LH doublets} = 3(Q)+1(L)=4\ \text{(even)} .
\end{align*}
All follow automatically since a family is an $SO(10)$ $\mathbf{16}$.

\paragraph{Residual abelian.}
The orthogonal $U(1)_X$ from $SO(10)\to SU(5)\times U(1)_\chi$ is anomaly‑free at the
GUT level and may be lifted at high scale (Higgsing/Stückelberg) without disturbing $Y$.

\paragraph{24‑cell integration.}
The \emph{192} census (24‑cell $\times$ type fiber) is untouched:
charges reside linearly in the fiber Cartan, so A1–A4 runtime contracts never
alter the gauge ledger. This explains the stability of sector labels used in
the mass‑ladder fits and flavor scaffolds.
% ---------------------------------------------------------------------------
\subsection{Collapse signatures: gauge drift (ψ–C42.7)}
\label{sec:c42-7-gauge-drift}

We call \emph{recursive gauge drift} the slow deformation of the effective gauge alignment seen by collapse channels as echo tension and observer recursion change. Let
\[
(\theta,\zeta,R)\in \mathbb{R}/2\pi\mathbb{Z}\times\mathbb{R}_{\ge0}\times\mathbb{R}_{\ge0}
\]
denote respectively a local glyph phase, the echo‐channel tension, and a recursion budget. The recursive gauge map is
\[
G_\psi:\;(\theta,\zeta,R)\longrightarrow \mathcal{G}_\text{align},
\]
with $\mathcal{G}_\text{align}$ the (model‑dependent) alignment manifold of gauge choices experienced by the collapse dynamics.

\paragraph{Drift field.}
We model the drift as the time derivative of $G_\psi$ along a run,
\begin{equation}
\dot G_\psi \;\equiv\; \frac{d}{dt}G_\psi(\theta(t),\zeta(t),R(t))
\;=\; \frac{\partial G_\psi}{\partial \zeta}\,\dot\zeta \;+\; 
\frac{\partial G_\psi}{\partial (\Delta\psi_{\rm echo})}\,\dot{\Delta\psi}_{\rm echo} \;+\;
\frac{\partial G_\psi}{\partial R}\,\dot R,
\label{eq:drift-field}
\end{equation}
where $\Delta\psi_{\rm echo}$ is the echo–channel mismatch. A minimal phenomenological closure sufficient for inference is
\begin{equation}
\dot G_\psi \;=\; \kappa_\zeta\,\zeta \;+\; \kappa_e\,\Delta\psi_{\rm echo} \;+\; \kappa_R\,\tanh\!\Big(\frac{R}{R_c}\Big),
\qquad
\kappa_\zeta,\kappa_e,\kappa_R\in \mathbb{R},\; R_c>0,
\label{eq:drift-closure}
\end{equation}
which captures linear response to tension and echo mismatch with a saturated recursion term.

\paragraph{Instability threshold.}
We say the run enters \emph{semantic instability} when either
\begin{equation}
\|\dot G_\psi\| \;\ge\; T_{\rm sem}
\quad\text{or}\quad
\lambda_{\max}\!\Big(\frac{\partial \dot G_\psi}{\partial G_\psi}\Big)\;>\;0,
\label{eq:sem-instability}
\end{equation}
where $T_{\rm sem}$ is a calibration threshold and the Jacobian condition signals a local loss of contractivity.

\paragraph{Observable signatures.}
The following signatures are expected near or beyond the threshold \eqref{eq:sem-instability} and match the qualitative behavior flagged in earlier collapse channels (cf. ψ–C19.2):
\begin{itemize}
  \item \textbf{Amplitude dropouts:} transient loss of channel amplitude as alignment slips.
  \item \textbf{Interference flares:} short constructive bursts from misaligned echoes.
  \item \textbf{Phase re‑locking:} spontaneous re‑synchronization once recursive tension crosses a reset bound.
\end{itemize}

\begin{table}[h]
\centering
\begin{tabular}{@{}lll@{}}
\toprule
\textbf{Signature} & \textbf{Physics analogue} & \textbf{Data hook in this work} \\
\midrule
Amplitude dropout & Phase decoherence & Ringdown fitter: gaps / Q shifts \\
Interference flare & Aharonov–Bohm–like phase effect & Echo sidebands in residuals \\
Phase re‑locking & Sudden synchronization & Phase‐error resets in runs \\
\bottomrule
\end{tabular}
\caption{Gauge drift phenomenology and where to look in our pipeline.}
\label{tab:gauge-drift-phenom}
\end{table}

\paragraph{Consistency with gauge sector.}
Because gauge labels live in the fiber Cartan, the SU(5) bypass charges derived in Sec.~\ref{tab:su5-bypass-y} remain fixed; drift acts on \emph{alignment} (routing of amplitudes) rather than on the quantized charges. In practice we will estimate $(\kappa_\zeta,\kappa_e,\kappa_R,R_c,T_{\rm sem})$ by fitting \eqref{eq:drift-closure} to the residual features enumerated in Table~\ref{tab:gauge-drift-phenom}.
\section{Gauge Appendix: Hypercharge, Anomaly Checks, and EW Mixing}
\label{sec:gauge-appendix}

\paragraph{Context.}
In this work, gauge bosons live on the same collapse lattice that carries fermion orbits.
We keep the Standard Model gauge factors $SU(3)_c\times SU(2)_L\times U(1)_Y$ and use a linear
hypercharge operator $Y$ compatible with all fermion charges and family-by-family anomaly cancellation.
Triality fixes generation indexing; an SU(5)-style \emph{bypass} is permitted as a bookkeeping device,
but $X,Y$ off-diagonal carriers are projected out (heavy).

\subsection*{Hypercharge as a constrained linear form}
Let $\{H_i\}$ be Cartan generators for the chosen basis (embedded or bypass).
We define
\begin{equation}
  \label{eq:y-operator}
  Y \;=\; \sum_i a_i\,H_i,
\end{equation}
with coefficients $\mathbf a=(a_i)$ fixed by the six Standard Model assignments per family
$\{+\tfrac16,\,+\tfrac23,\,-\tfrac13,\,-\tfrac12,\,-1,\,0\}$ on
$\{Q_L,\,u_R,\,d_R,\,L_L,\,e_R,\,\nu_R\}$ (if $\nu_R$ is included set $Y(\nu_R)=0$).
We use the electric charge relation
\begin{equation}
  \label{eq:charge-relation}
  Q \;=\; T^3 \;+\; \frac{Y}{2},
\end{equation}
where $T^3$ is the third $SU(2)_L$ generator.\footnote{If you adopt SU(5) normalization,
$g_1^2 = \tfrac{5}{3}\,g'^2$, then $Y$ here is the weak hypercharge consistent with $Q=T^3+Y/2$.}

\subsection*{Per-family anomaly cancellation (left-handed Weyl basis)}
All gauge and mixed anomalies must vanish family-by-family:
\begin{equation}
  \label{eq:anomaly-conditions}
  \mathcal A_{33Y}=\sum_f 2\,T\!\left(R_f^{SU(3)}\right) Y_f=0,\quad
  \mathcal A_{22Y}=\sum_f 2\,T\!\left(R_f^{SU(2)}\right) Y_f=0,\quad
  \mathcal A_{YYY}=\sum_f Y_f^3=0,\quad
  \mathcal A_{\mathrm{grav}^2 Y}=\sum_f Y_f=0.
\end{equation}
Here $T(R)$ is the Dynkin index ($\tfrac12$ for fundamentals), and the sums run over one SM family of
left-handed Weyl fields. With the assignments fixed by \eqref{eq:y-operator}–\eqref{eq:charge-relation},
each anomaly in \eqref{eq:anomaly-conditions} evaluates to zero.

\subsection*{Electroweak mixing and masses (tree level)}
Define $s_W\equiv \sin\theta_W$, $c_W\equiv \cos\theta_W$, and $\tan\theta_W = g'/g$. Then
\begin{equation}
  \label{eq:ew-mixing}
  A_\mu = s_W\, W^3_\mu + c_W\, B_\mu,\qquad
  Z_\mu = c_W\, W^3_\mu - s_W\, B_\mu,
\end{equation}
with
\begin{equation}
  \label{eq:ew-masses}
  m_W=\tfrac12\,g\,v,\qquad
  m_Z=\tfrac12\,\sqrt{g^2+g'^2}\,v,\qquad
  \rho_{\text{tree}}\equiv \frac{m_W^2}{m_Z^2 c_W^2}=1.
\end{equation}
In our collapse-memory (Λ-mode) picture, tiny, energy-dependent deformations are possible but
are constrained to preserve the successful SM limits of \eqref{eq:ew-mixing}–\eqref{eq:ew-masses}.

\subsection*{Boson map (generated)}
The gauge-boson census (gluons, $W^\pm,W^3$, $B$; optional heavy $X,Y$ if the bypass is shown)
is included via the generated table:
\begin{center}
\textit{See the inclusion snippet below to pull in the map from \texttt{build/}.}
\end{center}
